% (User input window)
% 2026-02-14 10:48:43（Asia/Singapore，UTC+08）

\paragraph{跨接口终端结论：碰撞压力、低温激发、顶端谱隔离与 Kronecker tick 的一维传输闭式}
本段继续把正文与附录中已建立的若干结构性断言，压缩为可独立引用的结论条目；
证明仅指向相应已完成的定理/命题/推论（或其可审计输出），因此不再重复细节推导。

\begin{conclusion}[碰撞压力的离散凸性与跨阶单调律：归一化压力斜率 \texorpdfstring{$g(q)$}{g(q)} 的单调上升与一致性审计]\label{con:xi-terminal-collision-pressure-discrete-convexity-crossq-g}
设
\(
S_q(m):=\sum_{x\in X_m}d_m(x)^q
\)
并以 \(r_q\) 表示其主增长常数（例如在 Perron--Frobenius 口径下 \(S_q(m)=\Theta(r_q^m)\)）。
则对任意固定 \(m\ge 2\)，函数 \(q\mapsto \log S_q(m)\) 在 \(q\ge 0\) 上为凸；
因此压力 \(P(q):=\log r_q\) 在整数点上满足离散凸性：对任意整数 \(q_0<q_1<q_2\)，令
\(
\theta=\frac{q_2-q_1}{q_2-q_0}
\),
则
\[
P(q_1)\le \theta P(q_0)+(1-\theta)P(q_2),
\qquad
\text{等价地 } r_{q_1}\le r_{q_0}^{\theta}r_{q_2}^{1-\theta}.
\]
同时，跨阶范数不等式导出：对任意 \(2\le p<q\)，在 \(|X_m|\sim\varphi^m\) 的尺度下有下界
\[
\log r_q \ \ge\ \frac{q}{p}\log r_p+\Bigl(1-\frac{q}{p}\Bigr)\log\varphi.
\]
等价地令
\(
g(q):=\frac{\log r_q-\log\varphi}{q}
\),
则 \(g(q)\) 在 \(q\ge 2\) 上单调不减，从而 \((r_q/\varphi)^{1/q}\) 随 \(q\) 单调上升；
该单调律可作为跨实现的回归一致性审计标量。
\end{conclusion}
\begin{proof}[证明]
见命题 \ref{prop:pom-pressure-convexity}（凸性与离散凸性）与命题 \ref{prop:pom-crossq-g-monotone}（跨阶下界与 \(g(q)\) 单调性）。证毕。
\end{proof}

\begin{conclusion}[低温端的严格稀疏激发律：\texorpdfstring{$\lambda(q)=2^q\bigl(1+(3/4)^q+O((5/8)^q)\bigr)$}{lambda(q)=2^q(1+(3/4)^q+O((5/8)^q))}]\label{con:xi-terminal-low-temperature-sparse-excitation-law}
令
\(
Z_L^{(q)}=\sum_{\boldsymbol{\ell}\in \mathcal C_L}\prod_i F_{\ell_i+2}^{\,q}
\)
并以 \(\lambda(q)\) 表示其长度增长常数（\(\lambda(q)=\rho(q)^{-1}\)，其中 \(\rho(q)\) 为方程 \(W_q(\rho)=1\) 的唯一正根；见定义 \ref{def:pom-multiplicity-real-q-pressure}）。
则 \(\lambda(q)\) 在 \(q\ge 0\) 上严格递增，\(\log\lambda(q)\) 严格凸且 \(\lambda(q)\) 为实解析函数；
并在低温极限 \(q\to\infty\) 给出精确渐近展开
\[
\lambda(q)=2^q\Bigl(1+\Bigl(\frac34\Bigr)^q+O\Bigl(\Bigl(\frac58\Bigr)^q\Bigr)\Bigr),
\qquad
\log_2\lambda(q)=q+\frac{1}{\ln 2}\Bigl(\frac34\Bigr)^q+O\Bigl(\Bigl(\frac58\Bigr)^q\Bigr).
\]
因此低温端由三常数 \(\{1,3/4,5/8\}\) 给出一条完全可审计的稀疏激发律：基态标尺为 \(2^q\)，首个激发态相对权重为 \((3/4)^q\)，而更高激发被 \((5/8)^q\) 量级统一压制。
\end{conclusion}
\begin{proof}[证明]
单调/凸性/解析性见定理 \ref{thm:pom-multiplicity-real-q-pressure}；
低温渐近展开见命题 \ref{prop:pom-multiplicity-lambdaq-large-q-asymptotic}；
将展开解释为低温激发律见命题 \ref{prop:pom-multiplicity-low-temperature-gibbs}。证毕。
\end{proof}

\begin{conclusion}[顶端谱的 Fibonacci 梯级缺口：第二、第三大纤维多重度闭式与奇偶极限常数]\label{con:xi-terminal-top-spectrum-fibonacci-ladder-gaps}
令 \(D_m:=\max_{x\in X_m}d_m(x)\)，并令 \(D_m^{(2)}\)、\(D_m^{(3)}\) 分别为第二、第三大纤维多重度（定义 \ref{def:pom-top-fiber-spectrum}）。
则当 \(k\ge 5\) 时存在闭式
\[
D_{2k}^{(2)}=F_{k+2}-F_{k-4},\qquad
D_{2k+1}^{(2)}=2F_{k+1}-F_{k-4},
\]
并且当 \(k\ge 6\) 时
\[
D_{2k}^{(3)}=F_{k+2}-F_{k-3}.
\]
奇数支在稳定相 \(m=2k+1\ge 19\) 进入精确闭式
\[
D_{2k+1}^{(3)}=2F_{k+1}-\bigl(F_{k-4}+F_{k-8}\bigr).
\]
从而顶端缺口在指数尺度上呈刚性 Fibonacci 梯级：对偶数支有
\(
D_{2k}-D_{2k}^{(2)}=F_{k-4}
\)、
\(
D_{2k}-D_{2k}^{(3)}=F_{k-3}
\)，并且奇数支的顶端两级缺口亦为显式 Fibonacci 组合。
进一步，顶端相对缺口
\(
G_m:=1-D_m^{(2)}/D_m
\)
具有奇偶极限常数
\[
\lim_{k\to\infty}G_{2k}=\varphi^{-6},\qquad
\lim_{k\to\infty}G_{2k+1}=\frac12\varphi^{-5},
\]
且伴随指数收敛。
\end{conclusion}
\begin{proof}[证明]
次大闭式见定理 \ref{thm:pom-second-max-fiber-closed-form}；
偶数支第三大闭式见定理 \ref{thm:pom-third-max-fiber-even-closed-form}；
奇数支第三大稳定闭式见注 \ref{rem:pom-third-max-fiber-odd-closed-form}；
梯级缺口汇总见注 \ref{rem:pom-top-fiber-spectrum-ladder}；
相对缺口极限常数见推论 \ref{cor:pom-top-gap-limit-constants}（其指数收敛由推论 \ref{cor:pom-second-max-gap-constant} 给出）。证毕。
\end{proof}

\begin{conclusion}[顶层谱的比例隔离带：除达到者外统一压至 \texorpdfstring{$\theta D_m$}{theta Dm}]\label{con:xi-terminal-top-spectrum-isolation-band}
令
\(
\mathcal{M}_m=\{x\in X_m:\ d_m(x)=D_m\}
\)
为最大达到者集合。
则对任意 \(x\notin\mathcal{M}_m\) 有 \(d_m(x)\le D_m^{(2)}\)，并且存在显式常数 \(\theta<1\)（可取全局 \(\theta=25/26\)，或渐近最坏常数 \(\theta_\star\approx 0.954915\)）使得对充分大的 \(m\) 有
\[
d_m(x)\le D_m^{(2)}\le \theta D_m\qquad(x\notin\mathcal{M}_m).
\]
因此顶端谱在比例意义下与近极值层之间存在稳定隔离带：极值层由有限相位承载，而全体非达到者在统一比例上被压入 \(\theta D_m\) 以下；
该隔离带为任何“以高阶矩强化尾部贡献”的推断提供无需额外假设的结构性输入。
\end{conclusion}
\begin{proof}[证明]
见注 \ref{rem:pom-maxfiber-isolation-band}（并结合推论 \ref{cor:pom-second-max-gap-constant} 对 \(\theta\) 的显式选取）。证毕。
\end{proof}

\begin{conclusion}[四相到二相的有效塌缩阈值：显式选择律 \texorpdfstring{$q_\delta$}{q_delta}]\label{con:xi-terminal-odd-fourphase-two-phase-collapse-threshold}
在奇数稳定相 \(m=2k+1\ge 13\) 的最大达到者族中，存在“偏置两相”与“平衡两相”的四相结构（注 \ref{rem:pom-max-fiber-achievers-bsplit}）。
记偏置最大分支
\(
A_k=F_{k+1}+F_{k-2}
\)，
平衡最大分支
\(
B_k=F_{k+1}
\)，
则比值 \(R_k=A_k/B_k\to 1+\varphi^{-3}\)。
对任意容差 \(\delta\in(0,1)\)，当
\[
q\ge q_\delta:=\Bigl\lceil \frac{\log(1/\delta)}{\log(1+\varphi^{-3})}\Bigr\rceil
\]
时，以 \(q\)-矩权重比较四相贡献，偏置两相的总占比至少为 \(1-\delta\)。
因此“最大达到者族在高阶矩下从四相有效塌缩到二相”可由单一闭式阈值 \(q_\delta\) 精确刻画。
\end{conclusion}
\begin{proof}[证明]
见命题 \ref{prop:pom-odd-fourphase-q-collapse-threshold}。证毕。
\end{proof}

\begin{conclusion}[时间纤维的立方几何压力：\texorpdfstring{$F_x(u)$}{F_x(u)} 的因子化与增长常数 \texorpdfstring{$r_\square(u)=\frac{1+\sqrt{5+4u}}{2}$}{r_square(u)=(1+sqrt(5+4u))/2}]\label{con:xi-terminal-timefiber-cubical-pressure-factorization}
令 \(C(n,k)\) 为 Fibonacci 立方 \(\Gamma_n\) 的 \(k\)-维子立方体数，并令 \(C_n(u)=\sum_{k\ge 0}C(n,k)u^k\) 为其 \(f\)-多项式（定义 \ref{def:pom-fibcube-fvector}）。
则 \(C(n,k)\) 满足二维递推
\[
C(n,k)=C(n-1,k)+C(n-2,k)+C(n-2,k-1),
\]
其生成函数为有理式
\[
\sum_{n\ge 0}C_n(u)z^n=\frac{1+z(1+u)}{1-z-(1+u)z^2},
\]
并具有显式系数公式
\[
C(n,k)=\sum_{j=k}^{\lceil n/2\rceil}\binom{n-j+1}{j}\binom{j}{k}.
\]
从而对每个固定 \(u\ge 0\)，指数增长常数存在且为
\[
\lim_{n\to\infty}\frac1n\log C_n(u)=\log r_\square(u),\qquad r_\square(u)=\frac{1+\sqrt{5+4u}}{2}.
\]
对一般纤维重构图 \(\mathcal F_x\)（分量长度多重集为 \((\ell_i)\)），其 \(f\)-多项式严格因子化
\[
F_x(u)=\prod_i C_{\ell_i}(u),
\]
从而“纤维立方压力” \(u\mapsto \log F_x(u)\) 完全由 \((\ell_i)\) 的加性组合确定，并在大分量极限中由同一闭式常数 \(r_\square(u)\) 统一控制。
\end{conclusion}
\begin{proof}[证明]
二维递推、生成函数与系数公式见定理 \ref{thm:pom-fibcube-fvector-closed}；
指数增长常数见推论 \ref{cor:pom-fibcube-fpoly-growth-constant}；
纤维 \(f\)-多项式因子化见推论 \ref{cor:pom-fiber-fpoly-factorization}。证毕。
\end{proof}

\begin{conclusion}[共振区的 Aitken 证书：\texorpdfstring{$(\sigma_{\mathrm{osc}},\delta_{\mathrm{ait}})$}{(sigma_osc,delta_ait)} 给出窗口长度的确定性选择律]\label{con:xi-terminal-resonance-aitken-certificate-window-law}
对共振区特征多项式 \(P_q(\lambda)\) 的根按模排序 \( |\lambda_{1,q}|>|\lambda_{2,q}|\ge|\lambda_{3,q}|\ge\cdots \)（其中 \(\lambda_{1,q}=r_q\)），定义
\[
\sigma_{\mathrm{osc}}(q):=\frac{|\lambda_{2,q}|}{r_q},\qquad
\delta_{\mathrm{ait}}(q):=\frac{|\lambda_{3,q}|}{r_q}.
\]
则比值估计
\(
R_m=S_q(m+1)/S_q(m)\to r_q
\)
的主误差在共振区表现为二周期交替几何项，衰减率由 \(\sigma_{\mathrm{osc}}(q)\) 控制；
Aitken \(\Delta^2\) 变换严格消去该交替主项，并将误差主控尺度降为
\(
\max\{\delta_{\mathrm{ait}}(q),\sigma_{\mathrm{osc}}(q)^2\}^m
\)。
因此给定目标相对误差阈值 \(\mathrm{tol}\in(0,1)\)，存在直接的窗口长度准则
\[
m\gtrsim \frac{\log(\mathrm{tol})}{\log \delta_{\mathrm{ait}}(q)}
\]
（在共振窗内可简化为 \(\delta_{\mathrm{ait}}\) 主导），从而“模 DP + 比值收敛 + Aitken”的端点估计具备完全由谱根指纹驱动的确定性证书。
\end{conclusion}
\begin{proof}[证明]
见推论 \ref{cor:pom-aitken-cancels-two-cycle}、定义 \ref{def:pom-resonance-aitken-fingerprints} 与注 \ref{rem:pom-aitken-window-audit}（含窗口长度准则与表 \ref{tab:fold_collision_resonance_aitken_certificate_q9_17}）。证毕。
\end{proof}

\begin{conclusion}[Kronecker 分母刻度的 \texorpdfstring{$W_1$}{W1} 精确闭式与单侧灵敏性：好侧二次、坏侧一次]\label{con:xi-terminal-kronecker-w1-denominator-one-sided-sensitivity}
设 \(p/q\) 为既约有理数，且 \(|\alpha-p/q|<1/q^2\)。
记 \(\delta=\alpha-p/q\)，令点集 \(P_q(\alpha)=\{\{t\alpha\}:t=0,\dots,q-1\}\) 的经验测度为 \(\mu_q\)，并记
\(
\mathsf{T}_q(\alpha)=W_1(\mu_q,\lambda)
\)。
则存在完全闭式的分母刻度公式：
\begin{enumerate}
  \item 若 \(\delta<0\)（坏侧，\(\alpha<p/q\)），则
  \[
  \mathsf{T}_q(\alpha)=\frac{1}{2q}+\frac{q-1}{2}\Bigl(\frac{p}{q}-\alpha\Bigr),
  \qquad
  \mathsf{T}_q(\alpha)=\frac12\,D_q^\ast(P_q(\alpha)).
  \]
  \item 若 \(\delta>0\)（好侧，\(\alpha>p/q\)），则
  \[
  \mathsf{T}_q(\alpha)=\frac{1}{2q}-\frac{q-1}{2}\Bigl(\alpha-\frac{p}{q}\Bigr)
  +\frac{q(q-1)(2q-1)}{6}\Bigl(\alpha-\frac{p}{q}\Bigr)^2,
  \]
  且该侧星差异退化为常数 \(D_q^\ast(P_q(\alpha))=1/q\)（每箱恰含一点）。
\end{enumerate}
因此同一“分母 tick”上，坏侧对 \(\delta\) 的响应为一阶线性（首箱双占/末箱空缺的结构性偏置），而好侧仅以二次项显影；
该单侧灵敏性在 \(W_1\) 口径下把“近似有理”的几何误差完全代数化，并给出可审计的确定性闭式预算。
\end{conclusion}
\begin{proof}[证明]
见定理 \ref{thm:kronecker-w1-denominator-closed-form} 与推论 \ref{cor:kronecker-w1-two-sided-interpretation}。证毕。
\end{proof}
